\documentclass[10pt]{article}
\usepackage[margin=20mm]{geometry}
\usepackage{amsmath,amssymb,mathtools,fontspec,unicode-math,xurl,hyperref}
\IfFontExistsTF{DejaVu Serif}{\setmainfont{DejaVu Serif}}{\setmainfont{Latin Modern Roman}}
\setmathfont{latinmodern-math.otf}
\setlength{\parindent}{0pt}\setlength{\parskip}{6pt}\sloppy
\title{Scientific review of
Oscillatory integral operators and variable Schrödinger propagators: beyond the universal estimates}
\author{Math 2025 — M25-0395 — targeted first-review recovery}
\date{28 October 2026}
\begin{document}
\maketitle
Reviewed work. Mingfeng Chen, Shengwen Gan, Shaoming Guo, Jonathan Hickman, Marina Iliopoulou and James Wright, arXiv:2407.06980v2, 26 October 2025, 84 physical PDF pages. The frozen annual entry is Geometric and Functional Analysis 35 (2025), 1425–1525, DOI https://doi.org/10.1007/s00039-025-00724-y. All page references below concern the exact v2 PDF, whose printed and physical page numbers agree. The journal version is not assumed to be identical to this manuscript.\par
\section*{Assessment and scope}
The manuscript develops a substantial route from analytic non-compression conditions to geometric maximal estimates and then to oscillatory integral and variable propagator estimates. The quantitative analytic sublevel arguments are an important part of that route. Nevertheless, the stated higher-dimensional Nikodym hypothesis admits a concrete five-dimensional phase that has precisely the compression the main results purport to exclude. This is a Major defect in Proposition 4.2(II) and Theorems 2.8(II) and 2.9(II), not merely an unverified difficulty in their proofs. Section 2 (\#sec:major) checks every hypothesis and supplies both geometric and oscillatory lower bounds.\par
Seven further findings concern explicitly identified definitions, normalizations, uniformity claims and local proof steps. Their proposed repairs have bounded scope. In particular, correcting these local steps does not repair the main counterexample. The assessment calls for a substantial revision of the high-dimensional Nikodym assertion. It does not assert failure of the three-dimensional Nikodym theorem, of the analytic sublevel propositions as standalone statements, or of every conditional estimate elsewhere in the paper.\par
All 84 pages were read, including the appendix and bibliography. The precise source and input manifests were checked before this work. No substantive old report was recoverable for this entry; the preserved sources and intake records were reused, while the present scientific analysis and report are new. The two primary references [KR, Guth] were also read in their entirety (eight and seven pages, respectively). This is an initial review and recovery deliverable. Subsequent independent scientific acceptance remains pending.\par
\section*{P1. A phase satisfying Hypothesis II with universal compression (Major)}
\subsection*{Location and exact claim}
Definition 2.7, pp. 10–11, permits arbitrary row and column sets \(\alpha,\beta\) of size \(d_{\rm crit}(n)\), with no requirement \(\alpha=\beta\). Proposition 4.2(II), proved on pp. 35–36, claims that Hypothesis II implies \(\mathrm{S}_{\mathrm N}(\phi;m_{\rm crit}(n);\kappa)\) for some \(\kappa>0\). Theorem 2.9(II), p. 11, then asserts a strict improvement over the universal Nikodym exponent for every \(p>p(n)\), and Theorem 2.8(II) asserts a corresponding strict propagator improvement. The following example satisfies the actual hypotheses in dimension five.\par
\subsection*{A real analytic Hörmander phase}
For \(x,y\in\mathbb R^4\), put \[B(t)=\begin{pmatrix}1-t^2&t\\t&0\end{pmatrix},\qquad
 H(t)=\operatorname{diag}(B(t),B(t)),\qquad
 \phi(x,t;y)=x\cdot y+\tfrac12 y^TH(t)y.\] This is a real polynomial on the required domain. We have \(\partial^2_{xy}\phi=I_4\), so H1 holds everywhere. Moreover \[B'(t)=\begin{pmatrix}-2t&1\\1&0\end{pmatrix},
 \qquad \det H'(t)=1.\] The Gauss normal has a nonzero last component: up to orientation it is \((-H'(t)y,1)/(1+|H'(t)y|^2)^{1/2}\). Consequently the Hessian in H2 is a nonzero scalar multiple of \(H'(t)\) and is nonsingular for every point. Thus H2 also holds everywhere. The inverse defining the curves is explicitly \[\Psi(\omega;t;y)=\omega-H(t)y,\] which is valid on the required small patch and extends polynomially beyond it. In particular there is no inverse-branch or domain obstruction to this example. Also \(\partial_y\phi(0,0;0)=0\).\par
Hypothesis II and the II parts of the main theorems do not impose the translation-invariant normalization \(\psi(0;y)=0\) of Definition 2.1. Our phase has \(\psi(0;y)=(y_1^2+y_3^2)/2\); this is allowed. It is the five-dimensional specialization of the dagger phase displayed in Section 3.4, p. 18, which is attributed there to Minxing Shen. The present review independently checks its compatibility with the new Hypothesis II; it does not claim authorship of that phase.\par
\subsection*{Verification of both parts of Hypothesis II}
Here \(d_{\rm crit}(5)=3\). Choose \[\alpha=\{1,2,4\},\qquad \beta=\{2,3,4\}.\] The relevant submatrix is \[[H(t)]_{\alpha,\beta}
 =\begin{pmatrix}t&0&0\\0&0&0\\0&t&0\end{pmatrix}.\] Its 19 nonempty square subminors, with row and column subsets as in \(\mathcal I(\alpha,\beta)\), consist only of \(0,t,t^2\). The nonzero span is exactly \(\operatorname{span}\{t,t^2\}\), independently of \(\omega,y\). Every Taylor polynomial of these functions of order at least two is the function itself. We may therefore take \(d=3\) in Definition 2.7(a). For every centre \(s\) and every \((\omega,y)\), the constant polynomial 1 is not in that span: evaluating a putative identity at \(t=0\) proves this immediately. Hence \(Z_{\rm N}\) is empty, as required.\par
For part (b), the infinite derivative-column matrix has constant rank \(r=2\) at every \(t\). The columns of \(t,t^2\) are independent because their first and second derivative rows have determinant 2. Equivalently, the normalized order-two jets of \(1,t,t^2\) at \(s\) form \[\begin{pmatrix}1&s&s^2\\0&1&2s\\0&0&1\end{pmatrix},\] whose determinant is 1. Thus both the finite-jet separation and the constant-rank condition hold uniformly over every parameter required by the definition. The calculation is unchanged if the full zero \(3\times3\) minor is omitted.\par
\subsection*{The sublevel conclusion is false}
In (5.5) take the admissible measurable maps \(\mathfrak D(\omega)=I_4\) and \(\mathfrak d(\omega)=\omega\); the latter maps the prescribed equal small balls to themselves. Then, up to sign, the determinant matrix in that equation is \[M(t)=I_4-H(t)
 =\operatorname{diag}\left(
 \begin{pmatrix}t^2&-t\\-t&1\end{pmatrix},
 \begin{pmatrix}t^2&-t\\-t&1\end{pmatrix}\right).\] Each block has rank one, so \(\operatorname{rank}M(t)=2\) for all \(t\). Every \(3\times3\) minor is identically zero, whereas the minor on rows and columns \(\{2,4\}\) equals 1. The sublevel set in (5.5) is therefore the whole time interval for every \(\sigma>0\). No positive \(\kappa\) can give the asserted bound \(C_\phi\sigma^\kappa\). This disproves Proposition 4.2(II), with no unverified parameter choices.\par
The decisive error is visible at the top of p. 36. The constant term of \(\det[H(t)\mathfrak D-I]_{\alpha,\beta}\) is written as \((-1)^{d_{\rm crit}}\). For a non-principal submatrix it is instead \((-1)^{d_{\rm crit}}\det[I]_{\alpha,\beta}\), which can vanish. In our example the whole left-hand side is zero; the displayed right-hand side would be \(-1\) plus a linear combination of \(t,t^2\), already impossible at \(t=0\). There is a second obstruction to that expansion for general \(\mathfrak D\): Cauchy–Binet introduces columns of \(H\) outside the chosen \(\beta\). Requiring a principal submatrix alone has therefore not been proved to repair the general argument.\par
\subsection*{Actual failure of the claimed maximal estimate}
For every small \(\omega\), select \(y=\omega\). The core curve is \(x=M(t)\omega\) and lies in the smooth three-dimensional set \[S=\{(x,t):x_1=-t x_2,\ x_3=-t x_4\}\] restricted to a fixed compact local patch containing these curves. Its \(\delta\)-neighbourhood has five-dimensional volume at most \(C\delta^2\). With \(g_\delta=\mathbf1_{N_\delta S}\), every selected horizontal \(\delta\)-tube is contained in \(N_\delta S\), so \(\mathcal N^\delta g_\delta(\omega)\geq c\) on a fixed positive-measure ball. Therefore \[\|\mathcal N^\delta\|_{L^p(\mathbb R^5)\to L^p(\mathbb R^4)}
 \gtrsim \delta^{-2/p}.\] This construction persists in arbitrarily small fixed source patches. For example \(p=4>p(5)=3\) gives the universal exponent \(1/2\). If the conclusion held with \(\beta<1/2\) for all \(\varepsilon>0\), choosing \(\varepsilon=(1/2-\beta)/2\) would contradict this lower bound as \(\delta\downarrow0\). Thus Theorem 2.9(II) itself fails.\par
\subsection*{Actual failure of the propagator improvement}
Choose a nonnegative smooth frequency cutoff \(\psi\), supported near zero and positive on a smaller ball. Choose a compact smooth amplitude equal to 1 on a sufficiently small product neighbourhood. Define \[\widehat f_\lambda(y)=e^{-i\lambda |y|^2/2}\psi(y).\] Stationary phase for the nonsingular quadratic form \(-|y|^2/2\) gives \(|f_\lambda(z)|\leq C\lambda^{-2}\) on \(|z|\leq C_0\lambda\); integration by parts outside that ball gives arbitrarily rapid decay. In particular \[\|f_\lambda\|_{L^q(\mathbb R^4)}\lesssim \lambda^{-2+4/q}.\] The output phase, writing \(\tau=t/\lambda\), is \(x\cdot y-\lambda y^TM(\tau)y/2\). For each block use the orthonormal basis \[v(\tau)=\frac{(1,\tau)}{\sqrt{1+\tau^2}},\qquad
 w(\tau)=\frac{(-\tau,1)}{\sqrt{1+\tau^2}} .\] The quadratic form vanishes in the two \(v\) directions and has eigenvalue \(1+\tau^2\) in the two \(w\) directions. Let \(x\) have \(w\) coordinates \(\lambda r_1,\lambda r_2\), with \(r_j,\tau\) in small fixed intervals. The two-dimensional stationary point in the \(w\) variables remains inside the region where the cutoffs are positive. Stationary phase in those variables gives a nonzero leading term of size \(\lambda^{-1}\); integration in the two free \(v\) variables has a positive leading coefficient. This remains bounded away from zero when the two physical \(v\) coordinates of \(x\) vary in sufficiently small fixed intervals: their additional Fourier factors are then uniformly close to 1. The resulting output region has volume comparable to \(\lambda^3\), since its free macroscopic coordinates are \(t,\lambda r_1,\lambda r_2\). All amplitude and stationary-point restrictions hold uniformly on this smaller patch. Consequently \[\|U^\lambda f_\lambda\|_{L^q(\mathbb R^5)}
 \gtrsim \lambda^{-1+3/q},\qquad
 \|U^\lambda\|_{L^q\to L^q}\gtrsim\lambda^{1-1/q}.\] For \(2<q<3=q(5)\) the latter is precisely \(\alpha_{\rm LS}(5;q)\). At \(q=5/2\) the input, output and ratio powers are respectively \(-2/5,1/5,3/5\). Taking a smaller positive \(\varepsilon\) than any claimed strict gain contradicts Definition 1.1 and Theorem 2.8(II). This proof only claims the lower bound on the valid local patch; it does not copy the overbroad assertion on all points of \(S^\lambda\) on p. 20 when a compact amplitude may vanish elsewhere.\par
\subsection*{Impact, repair boundary, and independent checks}
The counterexample reaches both high-dimensional II conclusions, not just an auxiliary determinant identity. It does not disprove Theorem 2.4(II) in dimension three, where the critical submatrix is the entire \(2\times2\) Hessian. Nor does it disprove Propositions 5.1 or 5.3 as independent analytic statements. Replacing Hypothesis II by a condition that genuinely implies the sublevel property in Definition 4.1 is necessary; assuming that property directly retains only the later conditional geometric argument, subject to the local corrections below. No sufficient weak replacement for Hypothesis II has been established in this review.\par
The retained script `independent\_phase\_checks.py` was written in this review and executed with exact rational polynomial arithmetic. It checks all 19 relevant minors, \(\det H'=1\), every one of the 16 \(3\times3\) minors of \(I-H\), its nonzero \(2\times2\) minor, normalized jets, the false constant term, and the rational exponents above. It executes no author code. The volume and stationary-phase arguments are the analytic proofs given here, not outputs claimed for that script.\par
\section*{P2. Hypothesis I includes the quantity it excludes (Minor)}
Location. Definition of \(\mathcal I(\alpha,\beta)\) and Definition 2.6, p. 10; compare the parenthetical explanation in the proof of Proposition 4.2(I), p. 35.\par
The defined set \(\mathcal I(\alpha,\beta)\) includes \((\alpha,\beta)\) itself. In the equation defining \(Z_{\rm K}\), set the coefficient of this full minor to 1 and all other coefficients to zero. The equation then holds for every \(t\) and every \(y\), for any phase. Thus \(Z_{\rm K}\) is the whole frequency ball and Hypothesis I is impossible. This is a direct logical obstruction, not a test on a special phase. The conditional higher-dimensional I statements are consequently vacuous as written; the issue is not a counterexample satisfying their impossible hypothesis.\par
The intended local correction is explicit in the proof on p. 35, which requires \(0<|\alpha'|=|\beta'|<d_{\rm crit}(n)\). Define separate index sets \(\mathcal I_{<}\) and \(\mathcal I_{\leq}\), use the proper-minor set for the Kakeya condition and retain the all-minor set where required in the Nikodym condition. Under Definition 2.1, the Hessian at \(t=0\) vanishes; evaluating at zero removes a possible constant term in a putative linear dependence. This verifies the local transition to Proposition 5.1 after the index correction. It does not validate a global change of \(\mathcal I\) in both hypotheses or repair P1. The error is classified as Minor because the specific intended index restriction is present in the proof and the local algebraic repair is exact.\par
\section*{P3. The inverse curves require a centering assumption (Minor)}
Location. Section 1.2, p. 4, immediately preceding (1.3). H1 alone does not ensure an inverse branch near \(\omega=t=y=0\) with \(x\) in the given unit ball. In dimension three consider \[\phi(x,t;y)=x\cdot y+t|y|^2+100y_1.\] This smooth real analytic phase satisfies H1 with mixed Hessian \(I_2\) and H2 with nonsingular Hessian \(2I_2\) up to the nonzero Gauss normal factor. Nevertheless \(\partial_{y_1}\phi=x_1+2ty_1+100>97\) throughout the stated domain. It cannot equal \(\omega_1\) for any \(\omega\) in a ball of radius \(\rho<1\) about zero. No reduction of that radius fixes the problem.\par
Specify a local base point and centre the \(\omega\) ball at the actual gradient, or first subtract the appropriate linear function of \(y\) so that \(\partial_y\phi(0,0;0)=0\). H1 then supplies the local inverse-function theorem after shrinking the patch. This subtraction does not change the mixed Hessian or curvature: for the oscillatory operator it is an input modulation, and for the propagator its Fourier modulation corresponds to translation of the input, both preserving the relevant norms. The geometric centre is relabelled accordingly. This is a necessary local normalization, not a counterexample to an estimate formulated with a correctly centred branch. The phase in P1 already has the correct centre.\par
\section*{P4. The Jacobian in Example 3.1 has the wrong middle sign (Minor)}
Location. Section 3.1, p. 16, the displayed formula for \(J\Phi\) and the following definitions of \(A,B,C\). Direct differentiation gives the spatial Jacobian \[D\omega(y)-\operatorname{diag}\left(t,\frac{t}{1-ty_2}\right).\] Its determinant is \[\frac{A(y)t^2-B(y)t+C(y)}{1-ty_2},\] with the same \(A,B,C\) printed in the source, rather than the displayed numerator \(At^2+Bt+C\). For the admissible smooth map \(\omega(y)=y\), at \(y=0,t=1/4\), the actual determinant is \(9/16\); the printed expression is \(25/16\). Both exact rational values were also checked by the independent arithmetic script.\par
The correction is simply the minus sign before \(Bt\). The example’s special choice \(A=B=C=0\) still makes the determinant zero, so the intended compression example is not refuted by this calculation. This finding concerns the general identity and its use for other choices of \(\omega\), not an unsupported claim that the entire example fails.\par
\section*{P5. Two semialgebraic proof steps need their missing justifications (Minor)}
Location. Section 4.1, Step 3, p. 26, and the multiplicity argument following (4.18), p. 29; Corollary A.5 and Lemma A.6, p. 82.\par
\subsection*{Compactness does not follow from semialgebraic selection}
Corollary A.5, reproduced from [Retained latex: \textbackslash{}cite[Lemma 2.2]\{KR\}], supplies a semialgebraic one-point-per-fibre section, not a compact section. The assertion on p. 26 that the selected \(\mathcal L_{\rm dir}\) and \(\mathcal L_{\rm ctr}\) are compact does not follow from it. For a decisive example to that inference, identify \(\mathbb R^4\) with \(\mathbb C^2\) and take the compact semialgebraic set \[A=\{(a,d):d=a^2,\ 1/4\leq |a|\leq1/2\}.\] A compact section over the annulus of \(d\) values would project continuously and bijectively onto that annulus. Compactness would make the inverse continuous, giving a continuous square root on an annulus around zero. On fundamental groups this would require \(2k=1\) for an integer \(k\), which is impossible. Thus even compact semialgebraic input does not in general admit a compact section. This refutes the invoked general inference; it is not presented as a fully realized counterexample to a particular tube family.\par
There is a local way to supply the needed parametrization: use a closed tube-containment set, taking a harmless slightly larger closed neighbourhood if necessary, and take the closure of the bounded semialgebraic section within that set. Semialgebraic closure preserves its dimension and bounded complexity. The projection still covers the desired base. Apply the compact parametrization lemma to this closure. Among the finitely many charts select one by the measure of its projection onto the relevant parameter base. This supplies the projected-measure lower bound actually needed later. Selecting merely a chart with large measure in the graph is not the same assertion. The later local inverse selection uses a nonsingular parameter projection on a chart; it does not require the closed graph to remain a one-point-per-fibre global section.\par
\subsection*{The inverse-composed map is not generally polynomial}
The maps \(\widetilde\Psi_{\iota,v}\) on p. 29 contain an inverse of a polynomial parameter map chosen earlier. Polynomiality is not preserved under this operation. Already \(\widetilde D(z)=z^2/2\), \(\widetilde P(z)=z/2\) on \(1/2<z<1\) and \(\widetilde\Psi(\omega,t;y)=\omega-ty\) give the component \[\widetilde P(\widetilde D^{-1}(y))-ty
 =\sqrt{y/2}-ty .\] The parameter Jacobian is nonzero on this patch, but the result is not a polynomial. Extra coordinates can be appended without removing this obstruction.\par
For the required Bézout bound, keep the original parameter \(z\) rather than substituting its inverse. If \(c\) coordinates of \(\widetilde D(z)\) (or \(\widetilde P(z)\)) are fixed at \(v\), impose those \(c\) polynomial equations together with the \(n-1\) polynomial output equations, with unknowns \(z\in\mathbb R^{n-1}\) and \(\eta\in\mathbb R^c\). This is a square polynomial system of bounded degree. At a point counted in (4.18), its Jacobian is, up to sign, the product of the nonsingular parameter Jacobian and the nonzero \((u,\eta)\) Jacobian already imposed there. The point is therefore an isolated complex zero. Bézout bounds the number of such zeros by a degree-dependent constant, irrespective of other positive-dimensional components. This restores the particular multiplicity estimate without asserting polynomiality of an inverse. These repairs address the two specified steps; they are not a new proof of every approximation and volume argument in Section 4.\par
\section*{P6. The multilinear estimate is not uniform in the first derivative alone (Minor)}
Location. Definition 6.10, p. 54, and the last quantitative sentence of Theorem 6.11, p. 55. The latter asserts dependence only on \(n,\varepsilon,M_1\) for curves transverse on a measurable set \(X\).\par
Take \(n=k=p=2\), and for an arbitrarily large integer \(N\) use the two-element smooth family \[\gamma_1(t)=0,\qquad \gamma_2(t)=\frac{\sin Nt}{N},
 \qquad -1<t<1.\] For each fixed \(N\) this is a uniform \(C^\infty\) family in the definition: every higher derivative has a finite bound, while \(M_1=1\) works for all \(N\). Choose a fixed \(0<c<10^{-2}\) and \(\delta=c/N\). Each of the two tube families consists of one \(\delta\)-tube. For zeros \(t_j=j\pi/N\) in \((-1/5,1/5)\) let \(X\) be the union of rectangles \[|x|<\delta/8,\qquad |t-t_j|<\delta/8 .\] There are comparable to \(N\) disjoint rectangles, each of area comparable to \(\delta^2\). They lie in both tubes and in \(B_{1/2}\). If a point of either core is within Euclidean distance \(\delta\) of such a rectangle point, its time differs from \(t_j\) by less than \(9\delta/8\). Hence the second curve’s slope has absolute value at least \(\cos(9c/8)>0.99\). Its unit tangent has wedge product greater than \(1/3\) with the vertical first tangent. The full infimum over nearby core points in Definition 6.10 is therefore bounded below by the fixed \(\nu=1/3\). Also \(\delta<\nu\) as required.\par
The left-hand side of (6.20) is at least \(c_1\delta\sqrt N\). With \(r=1/2\), its asserted right-hand side is \(C\nu^{-A}\delta^{1-\varepsilon}\), since both family cardinalities are one and the \(r\) exponent is zero. The ratio is bounded below by a constant times \(N^{1/2-\varepsilon}\), which diverges for \(0<\varepsilon<1/2\). This disproves the stated \(M_1\)-only uniformity, while respecting all smoothness and local transversality hypotheses. It does not disprove a bound with dependence on the higher derivatives of a fixed family.\par
The cited curved theorem [Retained latex: \textbackslash{}cite[Theorem 7, p.\textasciitilde{}6]\{Guth\}] assumes small slopes around a fixed coordinate axis for each family. Local transversality on \(X\) does not itself give that hypothesis. One repair is to allow dependence on appropriate higher derivative bounds and partition curves into uniformly controlled pieces with tangents in fixed small cones. For a fixed smooth phase such dependence is compatible with the manuscript’s phase-dependent constants. The example rules out retaining the stronger uniformity; it does not by itself refute the final phase-dependent estimates. A fully uniform first-derivative-only replacement is not supplied.\par
\section*{P7. Shrinking-cap symbol scales and packet centering (Minor)}
Location. Proposition 9.1, p. 73; the definition of \(f_T\), p. 74; Lemmas 9.2–9.3, p. 75.\par
The source places a smooth symbol on a cap of diameter \(O(R^{-1/2})\) but requires \(|\partial_y^\beta a_\theta|\lesssim R^{-|\beta|/2}\). These are not the bounds of the cap cutoffs already displayed correctly on p. 74. They would force such symbols to be rapidly small. Indeed, Taylor expansion along a coordinate from a point just outside the support, at distance \(O(R^{-1/2})\), gives \[\|a_\theta\|_\infty\leq C_L R^{-L}
 \quad\hbox{for every integer }L\geq1,\] since all jets vanish at the exterior point and the derivative of order \(L\) is \(O(R^{-L/2})\).\par
This contradicts the existence assertion of Proposition 9.1 for a fixed nonzero amplitude. Use the normalized constant Hörmander phase \(x\cdot y+t|y|^2\), a smooth amplitude equal to 1 near the origin, a fixed nonnegative nonzero \(f\), and \(\lambda=R\). At \(q=2\), the left side of (9.3) is bounded below on a fixed small ball for all large \(R\). On the right, every proposed cap symbol is \(O_L(R^{-L})\). The number of caps, the volume of \(B_{2R}\) and all prefactors grow only polynomially; the number of \(\gamma\) indices is fixed for fixed \(N,\varepsilon\). Taking \(L\) sufficiently large makes that entire term tend to zero. The \(R^{-N}\|f\|_2\) term also tends to zero. Thus the printed assertion cannot hold.\par
The needed cap scale is \(R^{|\beta|/2}\) (and \(\lambda^{|\beta|/2}\) in Lemma 9.3), as the neighbouring cutoff formula already shows. Separately, for a ball centred at a physical point \(z_B\), a coherent packet construction is to decompose \(h(y)=e^{i\phi^\lambda(z_B;y)}f(y)\) into \(h_{\theta,\omega}\) and define \[f_T(y)=e^{-i\phi^\lambda(z_B;y)}h_{\theta,\omega}(y).\] Then \(\sum_T f_T=f\), and the packet phase at \(z_B+z\) is \[\phi^\lambda(z_B+z;y)-\phi^\lambda(z_B;y)-\omega\cdot y,\] which is the phase compatible with the stated recentered tubes. The positive modulation printed for \(f_T\) on p. 74 has the opposite sign. When a unit-scale phase is used, its base point must also be \(z_B/\lambda\). The displayed reconstruction proves the corrected algebra and preserves packet norms.\par
These are local scale and modulation corrections. A unit-size symbol does not satisfy the printed negative derivative bounds, so it would be invalid to use one as a purported all-hypothesis counterexample to Lemma 9.2 itself. The contradiction above instead targets the existential Proposition 9.1, where it is valid. This review does not claim a new complete proof of the corrected reverse square function proposition.\par
\section*{P8. Local-to-global proof: Whitney separation and the summation range (Minor)}
Location. Proof of Lemma 9.5, pp. 79–80, particularly (9.21)–(9.24).\par
On p. 79, the relation \(Q_1\sim Q_2\) is defined by disjoint dyadic cubes with non-disjoint parents. Such siblings need not be separated: \([0,1)\) and \([1,2)\) have the same parent and distance zero. Products with common intervals give the same example in every dimension. Thus the stated inference \(\operatorname{dist}(Q_1,Q_2)\sim2^k\) is false, and the scale-decay factor in (9.21) cannot follow from that definition. Use the usual Whitney pairs of non-neighbouring cubes whose parents are neighbours. Their separation is comparable to their side length, and the complement of the diagonal admits the required bounded-overlap decomposition. This is the exact geometric property needed for the subsequent kernel bound.\par
In addition, the large-\(k\) series in (9.24) requires \(s\alpha<(n-1)/2\), not merely the lemma’s stated \(s\alpha<(n+1)/2\). For example, \(n=3,s=2,\alpha=1/2\) obeys the stated range and the available local \(L^\infty\to L^2\) estimate for the constant phase. Yet the large-\(k\) term is \(\min\{\sigma^{-2},2^k\}\), so its infinite sum diverges. This is a failure of that summation argument, not a counterexample to the lemma’s global conclusion.\par
For the application in odd dimensions \(n\geq3\), first separate the range \(\alpha\geq\alpha_{\rm H}(n,s)\): its requested global exponents are already supplied by the universal theorem. In the remaining range, using \[\alpha_{\rm H}(n,s)=\frac{n+1}{2s}-\frac{n-1}{4},
 \qquad 2\leq s\leq q(n),\] gives \(s\alpha<s\alpha_{\rm H}(n,s)\leq1\leq(n-1)/2\). The strict summability inequality then holds, and the geometric series can be evaluated as intended. At the boundary \(n=3,s=2\), the case of equality is assigned to the universal estimate. Together with the corrected Whitney decomposition this supplies the missing local proof organization. No false global theorem is alleged on the basis of the divergent intermediate bound.\par
\section*{Coverage, counter-readings, and limits of the conclusions}
The following record separates reading and checked deductions from stronger conclusions that are not warranted.\par
Source pages | Scope and result of the audit\par
1–13 | Definitions, local domains, all quantifiers in Hypotheses I/II, three- and higher-dimensional theorems, universal exponents, and Proposition 2.10. P1 and P2 concern genuinely different conditions; P3 records the inverse-branch normalization.\par
14–20 | Examples 3.1–3.3, contact-order comparisons, Pierce–Yung application and sharpness constructions. The dagger phase is credited to the source; its additional compatibility with Hypothesis II is independently established in P1. P4 corrects the general Jacobian without invalidating its zero example.\par
21–32 | Grains, tube selection, sublevel-to-non-concentration estimates, Taylor approximation, parametrization, area formula and multiplicity control. P5 identifies the two missing local justifications. The Nikodym sublevel input must have the supremum in the centre variable and measure in time; the transposed wording in (4.16)(II), p. 28, should be aligned with the actual use later.\par
33–46 | Propositions 5.1 and 5.3, finite jets, Wronskians, analytic exceptional sets, uniform coefficient separation and compactness. These arguments were checked as analytic statements under their own hypotheses. Their application on p. 36 fails for the algebraic reason in P1. In the Taylor-remainder estimate for Proposition 5.3, retain the factor \(1+\|\mathfrak p\|\) that is present in the final bound, rather than bounding the remainder of \(f\) by \(\|\mathfrak p\|\) alone.\par
47–55 | Multiplicity duality, linear-to-multilinear reduction, transversality, exponent ranges and the curved universal input. P6 isolates the false stronger dependence on \(M_1\). The phase-dependent use can allow higher derivative bounds.\par
56–70 | Polynomial partitioning, cell/algebraic alternatives, dimension descent, stopping scales, mass allocation, dyadic losses and non-concentration gain. The cell statements are used inside the prescribed ball; an unrestricted assertion that every polynomial complement cell is bounded would be false. Harmless enlargements of the loss constants allow the gain parameter to be chosen strictly between zero and one. No separate counterexample to the resulting conditional estimate was established.\par
71–72 | Assembly of linear and multilinear estimates from the sublevel hypotheses. The Nikodym route uses the general centre-separated statement; a superfluous translation-invariance restriction in its intermediate wording should not be used to conceal the P1 failure.\par
72–81 | Universal estimates, reverse square function, packets, localized kernels, maximal-function transfer, local-to-global step and the final exponent choices. P7 and P8 specify local corrections. The full high-dimensional II result remains contradicted by P1 regardless of them.\par
81–84 | Real algebraic tools, appendix hypotheses, acknowledgements and complete reference list. Primary versions of [Retained latex: \textbackslash{}cite] [citation keys: KR; Guth] were separately acquired and read; this does not mean every cited book or external proof was newly read.\par
Several tempting overstatements were explicitly rejected. The extra constant quadratic term in the dagger phase does not exclude it from Hypothesis II; Definition 2.1 is not imposed there. Conversely, the same example does not establish a failure in dimension three. Failure of the \(M_1\)-only bound does not by itself disprove a fixed-phase estimate with higher-derivative constants. An inverse-composed parametrization need not be polynomial, but Bézout can still be applied in the original variables. The local divergence in the proof of Lemma 9.5 does not imply that its global conclusion is false. The source explicitly omits some details of the positive-definite corollary, Theorem 2.5; this report does not replace that omission with a claim of a new complete proof.\par
The corrected proper-minor definition, centering operation, Jacobian sign, original-variable multiplicity argument, packet reconstruction and summation range split have the specific verification described above. A new high-dimensional non-compression theorem, the complete corrected reverse square function argument, and a first-derivative-only curved theorem have not been proved here. These limits are material and are not erased by the delivery status.\par
\section*{Supporting checks}
The exact main source contains 863312 bytes and has SHA256 \[\text{\small\texttt{ b144c8027f924075c82f5610af31f7202b03d4e57fb344e94686e3fcc2372c24}}.\] Its preserved manifest and byte-matched copies retain their original execution-path attribution. The available arXiv landing page independently identifies the six authors and v2 date. A bounded search did not identify a necessary separately distributed supplement or author code for the reviewed proof; absence of code is not treated as a mathematical defect. No author code was executed.\par
The per-paper evidence includes the source text extraction, the original intake bindings, both exact reference PDFs and their retrieval and extraction records, the exact-arithmetic script and execution records, and the source-image inspection record. Eighteen main-source pages were actually displayed and observed: 1, 4, 10, 11, 16, 18–20, 26, 29, 35–36, 55, 73–75, 79 and 82. The two additionally observed reference pages were [KR, p. 3] and [Guth, p. 6]. This is separate from the complete textual reading of all 84 main-source pages and all 15 pages of those two references. An initially truncated image response was not counted as inspection; its pages were subsequently displayed individually.\par
All supporting scripts, compilation cache, logs, historical attempts, current TeX, actual final PDF, parse receipts and actual report-page inspections are retained under this paper’s assigned directories and indexed by path, bytes and SHA256 in its support manifest. A successful final build is credited only to the exact current TeX, compiler, source and unchanged per-paper dependency cache. Rasterization alone is not recorded as visual inspection.\par
The controller’s observation of the authorized Academic Researchers Chat6Pro interface is recorded separately; backend equivalence was not independently verified here. The scientific reasoning, exact calculation, writing, compilation and image observation were performed by the root scientific session. No independent acceptance is claimed; the scientific decision remains `pending` and the human review list remains empty.\par
9 N. H. Katz and K. M. Rogers, On the polynomial Wolff axioms, arXiv:1802.09094v3, 22 July 2018, 8 pp. https://arxiv.org/pdf/1802.09094v3. The complete acquired version was read; Lemma 2.2 and the parametrization argument are the relevant dependencies. L. Guth, A short proof of the multilinear Kakeya inequality, arXiv:1409.4683v2, 5 November 2014, 7 pp. https://arxiv.org/pdf/1409.4683v2. The complete acquired version was read; Theorem 7, p. 6, is the relevant curved estimate.\par
\end{document}
